\begin{claim}{D3.1}{Proved}{The obstruction map and the Kuranishi germ}
The primary Kuranishi bracket is zero, the Kuranishi germ of $X_c$ is
a smooth one-dimensional germ, and the $c$-family is locally
semiuniversal. The equality $h^2(T_X)=h^1(T_X)=1$ does not identify the
nonzero obstruction space with a nonzero obstruction map.
\end{claim}
\begin{proof}
Let $\alpha$ be the generator in D2.1. The constructed holomorphic
family gives an actual convergent deformation tangent to $\alpha$.
Differentiating its integrability equation to second order gives
$\bar\partial\alpha_2=-[\alpha,\alpha]/2$, so the cohomology class
of $[\alpha,\alpha]$ is zero. Bilinearity proves that the bracket
$H^1(T_X)\times H^1(T_X)\to H^2(T_X)$ vanishes on every pair.
This argument alone would not exclude higher obstructions, so we
make the analytic argument as well. A Kuranishi germ is a closed
analytic subgerm of its tangent space $H^1(T_X)=\C$. Its classifying
map from the actual parameter disc has nonzero derivative by D2.1.
Composing with this embedding gives a holomorphic map of discs with
nonzero derivative; the inverse function theorem makes its image
contain a neighbourhood of zero in $\C$. Every equation of the
Kuranishi germ vanishes on that neighbourhood and is therefore zero.
The germ, including its analytic structure, is the smooth line.
The same inverse-function argument makes the classifying map an
isomorphism of germs. Pulling back the semiuniversal deformation
therefore realizes every sufficiently small deformation in the
$c$-family. This is semiuniversality; the positive-dimensional
automorphism group is not being discarded or interpreted as a
fine-moduli assertion.
\end{proof}
